\section{Giriş}
\label{CH:Giris}

\subsection{Problem ve motivasyon}

Türk İşaret Dili, işitme engelli bireylerin hem kendi aralarında hem de çevreleriyle
iletişim kurarken kullandığı doğal dildir~\cite{TIF2023}. İşaret dilini bilmeyen
kişilerle kurulan iletişim ise çoğu zaman yazı yazmaya ya da bir yakının
tercümanlığına bağlı kalmaktadır. Hastane, banka, okul ve resmî daire gibi ortamlarda
bu durum hem zaman kaybına hem de mahremiyet sorununa yol açmaktadır.

Projenin çıkış noktası, bu iletişim boşluğunda en yaygın cihazı, yani akıllı telefonu
kullanmaktır. Eldiven ya da derinlik kamerası gibi ek donanım gerektirmeyen bir çözüm
hem maliyeti düşürür hem de günlük kullanımda kabul edilebilirliği artırır. Sistemin
internet bağlantısına ihtiyaç duymaması, şebekenin zayıf olduğu yerlerde de
kullanılabilmesini ve kamera görüntüsünün cihazdan çıkmamasını sağlar.

\subsection{Amaç ve kapsam}

Projenin amacı, telefon kamerasının önünde yapılan bir TİD kelimesini tanıyıp yazıya
çeviren ve tamamen cihaz üzerinde çalışan bir Android uygulaması geliştirmektir.
Kapsam şu şekilde sınırlandırılmıştır:

\begin{itemize}
  \item Tanıma kelime düzeyindedir (izole işaret). Sözlük, AUTSL veri setindeki 226
        işaretle sınırlıdır.
  \item Art arda yapılan işaretler sürekli modda tek tek tanınıp yan yana yazılır;
        bunların dilbilgisel Türkçe cümleye çevrilmesi kapsam dışıdır.
  \item Parmak alfabesi (harf harf heceleme) veri setinde bulunmadığı için desteklenmez.
\end{itemize}

\subsection{Başarı ölçütleri}

Tez önerisinde belirlenen ölçütler Tablo~\ref{tab:olcutler} içinde verilmiştir. Bu
ölçütlerin hangilerinin sağlandığı Bölüm~\ref{CH:Degerlendirme} içinde
tartışılmaktadır.

\begin{table}[H]
\centering
\caption{Tez önerisindeki değerlendirme ölçütleri ve hedefler}
\label{tab:olcutler}
\begin{tabular}{@{}lll@{}}
\toprule
\textbf{Ölçüt} & \textbf{Ölçüm ortamı} & \textbf{Hedef} \\
\midrule
Top-1 doğruluk & AUTSL test kümesi & $\geq$ yüzde 85 \\
Gerçek ortam doğruluğu & Gönüllü testi (30 işaret) & $\geq$ yüzde 75 \\
Model boyutu & \texttt{.tflite} dosyası & $\leq$ 10 MB \\
İşleme hızı & Orta segment Android telefon & $\geq$ 15 kare/s \\
İnternet bağımlılığı & Uçak modu & Tam çalışma \\
\bottomrule
\end{tabular}
\end{table}

\subsection{Ekip ve iş bölümü}

Proje üç kişilik bir ekip tarafından yürütülmüştür. Her alanın bir sorumlusu
belirlenmiş, ancak veri, model ve uygulama aynı hattın parçaları olduğu için kararlar
birlikte alınmıştır: veri hazırlığı, veri hattı ve normalizasyon (Mustafa Küpeli);
yapay zekâ modeli, eğitim ve optimizasyon (Mustafa Göktepe); mobil uygulama, kullanıcı
arayüzü ve donanım entegrasyonu (Ahmet Çelik). Çalışma ortak bir Git deposunda
yürütülmüş, her deney \texttt{experiments/} klasöründe tarih, tohum ve sonuçlarıyla
kaydedilmiştir.

\subsection{Tezin düzeni}

Bölüm~\ref{CH:Literatur} işaret dili tanıma alanındaki çalışmaları ve Türk İşaret Dili
veri setlerini özetlemektedir. Bölüm~\ref{CH:Veri} veri setini ve ön işlemeyi,
Bölüm~\ref{CH:Model} model mimarilerini, eğitimi ve hata analizini anlatmaktadır.
Bölüm~\ref{CH:Uygulama} modelin telefona aktarılmasını ve mobil uygulamayı,
Bölüm~\ref{CH:Degerlendirme} başarı ölçütlerine göre değerlendirmeyi içermektedir.
Bölüm~\ref{CH:Sonuc} sonuçları ve gelecek çalışmaları özetlemektedir.
